
de @@M0@@
dans @@M1@@, supposé partout de rang @@M2@@. Dans loc. cit., cette formule résulte d'un calcul très simple de faisceaux @@M3@@. Grâce à ce qu'on a vu dans 4.2, la formule précédente implique la même formule dans @@M4@@, mais seulement \uline{modulo torsion}. Il est assez extraordinaire que la formule (4.6) dans @@M5@@ lui-même ne soit pas prouvée à l'heure actuelle\SGAAuthorFootnote{*}{cf. note de la page 23.\newline (**)}, ni même le cas particulier obtenu en faisant @@M6@@, donnant la self-intersection de @@M7@@:
@@M8@@
même si @@M9@@ est le corps des complexes !

On se heurte à un problème analogue lorsqu'on tente de calculer l'anneau de Chow de la variété éclatée @@M10@@ de @@M11@@ le long de @@M12@@, en calquant la méthode suivie pour le calcul de @@M13@@ [Exp. VII 4.8], et qui marche également (sans tensoriser par @@M14@@) pour le calcul de @@M15@@. La formule-clef pour cette détermination (avec les notations de loc. cit.) est
@@M16@@
formule établie à l'heure actuelle dans @@M17@@ (pour @@M18@@) grâce à un calcul de @@M19@@ (Exp. VII 3.4), et dont la validité dans @@M20@@ (pour @@M21@@) resterait à établir.

Signalons que même les analogues @@M22@@-adiques des formules (4.6), (4.7), (4.8) ne sont pas prouvés à l'heure actuelle; lorsque @@M23@@ est la classe de cohomologie d'un cycle algébrique, (4.6) et (4.8) sont vraies en tous cas mod. torsion.

\textbf{4.5.} Dans cette section et les suivantes, nous allons donner des exemples où dans le diagramme (4.1) l'homomorphisme @@M24@@ respectivement l'homomorphisme @@M25@@ n'est pas un isomorphisme. Il revient évidemment au même, pour @@M26@@ donné, de dire que
@@M27@@
est un isomorphisme, ou un monomorphisme, ou un épimorphisme. D'ailleurs l'image de @@M28@@ est le sous-anneau de @@M29@@ engendré par les classes de Chern des faisceaux localement libres sur @@M30@@. Dans le cas où @@M31@@ est engendré comme algèbre sur @@M32@@ par les classes des faisceaux inversibles sur @@M33@@, on voit tout de suite que le sous-anneau envisagé de @@M34@@ est aussi le sous-anneau engendré par @@M35@@, la composante de degré @@M36@@ de @@M37@@.
